\section{Source-Frame and Functional-Outcome Applicability}\label{sec:r41-utility}
\paragraph{Selection consequence} All reporting statistics concern mitigation results in studies admitted through C1--C5, including an empirical inference evaluation. They are not a census of all defense research. SafeKV, CachePrune and OTRO also demonstrate why this is not a claim that every included study is attack-first. Adjacent defense-oriented work remains outside the denominator when qualification is unresolved or a required criterion is not established.

\paragraph{Revised R8 rule} A quality measurement is applicable when the indexed intervention can change prompts, model content, sampling, approximate computation or component coverage. A lossless wrapper, fixed-trace transformation or exact logical computation can instead support a source-scoped preservation claim. Such a record is exempt from the content/function-quality denominator; its timing, resource use and end-to-end task outcomes still require their own evidence. Padding an image's patch grid and inserting null prompt tokens are not automatically content-preserving: they change model processing, unlike wire padding.

The 31 records comprise 11 potentially quality-changing interventions, 18 scoped preservation witnesses and two standalone cost/calibration units without an output-quality comparison. Projection gives eight quality-applicable families: three B, one P and four U; eleven families are exempt from this narrow query. Thus the old 3/16 and 5/16 counts are superseded by 3/8 and 4/8 for functional quality. These describe the selected source frame and the named outcome. The old source cells remain archived; \path{data/r41-utility-scopes.json} provides every new applicability judgment, original status, locator and reason.

\paragraph{Agent subset} The map-selected four-family inventory contains six records. UI timestamp transformations and Dangol capture/live transport preserve saved actions/messages for the narrow functional object. WRA returns only the real output despite parallel decoy execution, so it also supports a narrow real-answer-content preservation claim. The quality-applicable agent-associated subset is therefore the cloud benchmark alone (0/1 B, 1/1 B+P family; 0/2 B, 1/2 B+P records). Its prompt/temperature conditions measure fingerprinting, not MARBLE task completion. Added execution is a separate outcome: WRA supplies no matched end-to-end task score for the decoy condition (0/1), and no explicit same-condition task score is linked for any of the four mapped families (0/4). The latter is a broad inventory, not four execution-changing cover implementations.

\paragraph{Positive functional routes} The measured-quality joins remain CachePrune's two target outcomes and SparLeak's intensity-specific quality result. Separately, SPEC's payload framing and Dangol's capture/live records combine scoped preservation with bound inference/cost objects. These records support their stated payload, trace or message claims, rather than an unmeasured task/deadline guarantee. Timing and placement limitations of their costs remain unchanged.

\begin{table}[ht]\centering\small
\caption{Functional-quality denominator after applicability and source-removal checks. Preservation and component-only units are exempt.}\label{tab:r41-removal}
\begin{tabularx}{\linewidth}{@{}Yrrr@{}}\toprule
Source inventory & Families & Quality B/applicable & Quality B+P/applicable\\\midrule
All source families & 19 & 3/8 & 4/8\\
Without CachePrune & 18 & 2/7 & 3/7\\
Without final-search records & 15 & 1/5 & 2/5\\
\bottomrule\end{tabularx}\end{table}

\begingroup\footnotesize\setlength{\tabcolsep}{3pt}
\begin{longtable}{@{}p{.23\textwidth}p{.12\textwidth}p{.08\textwidth}p{.18\textwidth}p{.33\textwidth}@{}}
\caption{Complete R8 applicability overlay for all 31 Type I/II records. Relevant = functional quality applies; preserved = scoped content/action/function witness; component = no output-quality unit. Archived source statuses remain available.}\label{tab:r41-utility-records}\\
\toprule Record & Scope & Old/new R8 & Source locator & Outcome-specific reason\\\midrule\endfirsthead
\toprule Record & Scope & Old/new R8 & Source locator & Outcome-specific reason\\\midrule\endhead
CACHEPRUNE-\allowbreak{}context~\cite{wu2026cacheprune} & Relevant & B/B & 2.2.2; 4.2; 5.3; 5.6 & Position-independent KV reuse approximates attention dependencies; generation quality depends on recomputation.\\
CACHEPRUNE-\allowbreak{}selective~\cite{wu2026cacheprune} & Relevant & B/B & 2.2.2; 4.2; 5.3; 5.6 & Same approximate segment-reuse configuration as contextual recovery; linked F1/ROUGE results apply.\\
CARLINI-\allowbreak{}fixed-\allowbreak{}rate~\cite{carlini2024remote} & Preserved & U/-- & VII; Figures11-12 & Fixed-rate delivery queues existing tokens and emits empty cover messages; narrow token-content preservation excludes delay and live-placement claims.\\
DANGOL-\allowbreak{}capture~\cite{dangol2026workflow} & Preserved & U/-- & 9.8-9.9; Table6 & Shaping transforms a saved capture with fixed application messages; no new output-quality experiment is needed for those fixed messages.\\
DANGOL-\allowbreak{}livebinding~\cite{dangol2026workflow} & Preserved & U/-- & 9.7; OHTTP paragraph after Table5 & Encapsulation and traffic cadence preserve application message bodies; task deadlines and end-to-end workflow completion remain separate outcomes.\\
EARLY-\allowbreak{}K-\allowbreak{}calibration-\allowbreak{}oracle~\cite{earlybird2025} & Preserved & U/-- & V-A; TableVI & Prefix-sharing granularity and oracle calibration do not edit model input or the logical prefix KV computation; this is narrower than an online function measurement.\\
EARLY-\allowbreak{}semantic-\allowbreak{}deidentification~\cite{earlybird2025} & Relevant & U/U & V-B; Figure12 & PII is replaced before semantic-cache processing and restored afterwards; content-sensitive processing needs functional evidence.\\
GAO-\allowbreak{}zero-\allowbreak{}copy-\allowbreak{}off~\cite{gao2025iknow} & Preserved & --/-- & 7.1 & Disabling zero-copy changes loading and memory sharing while retaining the logical model; loading/memory costs remain applicable.\\
KVGOV-\allowbreak{}HMAC-\allowbreak{}cost~\cite{addagada2026kvgov} & Component & U/-- & 5.9 & Key-derivation microcost has no model-output/function comparison unit; its resource cost is retained.\\
KVGOV-\allowbreak{}hardware-\allowbreak{}gap~\cite{addagada2026kvgov} & Component & U/-- & 5.7-5.8 & Unprotected TTFT calibration is not a modified model-output condition; its timing result is not a quality measurement.\\
KVGOV-\allowbreak{}partialflush-\allowbreak{}sim~\cite{addagada2026kvgov} & Preserved & U/-- & 4.1-4.2; 5.1-5.3 & Namespace salting and cache reuse/eviction policy change cache visibility and recomputation, not prompt content or the logical prefix computation; scope is the simulated policy.\\
NET-\allowbreak{}dummy-\allowbreak{}padding-\allowbreak{}simulation~\cite{netecho2025} & Preserved & U/-- & VII; TableIX & Synthetic packet insertion and padding operate on saved traces, leaving their stored application content fixed.\\
NET-\allowbreak{}upstream-\allowbreak{}local~\cite{netecho2025} & Preserved & U/-- & IV-B; IV-C-E; TablesIII-IV; Figure5 & Padding and batching alter framing/grouping of the token stream. The exemption concerns that content projection, not every frontend or online deployment behavior.\\
OTRO-\allowbreak{}epochpad~\cite{lee2026otro} & Preserved & --/-- & V-C-D; VI; VII-C-E & Dummy memory accesses and internal scheduling preserve the logical token sequence; padded execution cost is still distinct from unpadded cost.\\
OTRO-\allowbreak{}unpadded~\cite{lee2026otro} & Preserved & --/-- & V-C; VII-C-D & Internal oblivious-tokenization scheduling preserves token IDs relative to the non-ORAM baseline; only its own measured costs may be reused.\\
SAFE-\allowbreak{}post-\allowbreak{}trigger~\cite{chu2025safekv} & Preserved & U/-- & 4.2.3; 4.3.1; 7.2-7.4 & Visibility bits and exact-prefix cache lookup govern reuse without rewriting prompts or approximate KV substitution; this exempts narrow generation quality, not service performance.\\
SHAPE-\allowbreak{}padding-\allowbreak{}estimate~\cite{hadad2026shape} & Relevant & U/U & IX-B1 & Forces a worst-case image grid/patch configuration; unlike wire padding this changes model preprocessing, so quality preservation is not automatic.\\
SPARLEAK-\allowbreak{}replacement~\cite{chen2026sparleak} & Relevant & B/B & 7.5; Table4 & Random token replacement changes generated content; source reports intensity-matched ROUGE.\\
SPEC-\allowbreak{}padding-\allowbreak{}REST~\cite{wei2024speculation} & Preserved & --/-- & 6.1; Figures6-7; Table2 & Padding/aggregation wraps already generated token payloads; live/wire placement and interactive delay remain separately scoped.\\
SPLICE-\allowbreak{}CTBF-\allowbreak{}signal~\cite{spliceleak2026} & Relevant & U/U & VI-B; VII-C; TableII & Changes selective boundary recomputation in non-prefix approximate fusion; unchanged generation quality is not established by a fixed FLOP budget.\\
SPLICE-\allowbreak{}QCP-\allowbreak{}cost~\cite{spliceleak2026} & Relevant & U/U & VI-A; VII-C; TableII & Appends null tokens to a private prefix; no output-invariance or attention-mask witness is bound in the inspected scope.\\
TDXRAY-\allowbreak{}oblivious-\allowbreak{}map~\cite{hornetz2026tdxray} & Relevant & B/B & 6.1; Figure8; Table3 & Truncated token-frequency/vocabulary mapping has a component coverage outcome; it remains a functional-quality query.\\
TIME-\allowbreak{}length-\allowbreak{}instruction~\cite{timewilltell2024} & Relevant & U/U & VII-B & A fixed-length explanation instruction changes the requested answer, unlike transport padding.\\
TIME-\allowbreak{}padding-\allowbreak{}estimate~\cite{timewilltell2024} & Preserved & U/-- & VII-C & System-level padding delays and frames existing output tokens; it is separate from the answer-length instruction.\\
TRAFFIC-\allowbreak{}prompt-\allowbreak{}reformulation~\cite{pouryousef2026traffic} & Relevant & P/P & V-C1; Figure9a & Prompt wrapping changes model input. Semantic preservation is asserted rather than matched task-scored.\\
TRAFFIC-\allowbreak{}temperature~\cite{pouryousef2026traffic} & Relevant & U/U & V-C1; Figure9a & Temperature changes decoding and can change generated content.\\
UI-\allowbreak{}delay-\allowbreak{}offline~\cite{uitraces2026} & Preserved & U/-- & 6.1; Figure5; AppendixA.5 & Offline timestamp edits retain saved action types and structure; the exemption does not establish live task or timeout preservation.\\
WHISPER-\allowbreak{}batching-\allowbreak{}replay~\cite{whisperleak2025} & Preserved & U/-- & 6.2; Figure7 & Grouping stored packet features leaves saved model text fixed; streaming response time remains an unmeasured deployment outcome.\\
WHISPER-\allowbreak{}injection-\allowbreak{}replay~\cite{whisperleak2025} & Preserved & U/-- & 6.3; Table5 & Synthetic packet injection is a trace transformation rather than model-content editing.\\
WHISPER-\allowbreak{}randompadding-\allowbreak{}capture~\cite{whisperleak2025} & Preserved & U/-- & 5; 6.1; Tables3-4 & Additional random fields change serialized event length, not decoded token content.\\
WRA-\allowbreak{}decoy-\allowbreak{}execution~\cite{jeong2026wra} & Preserved & P/-- & 8.1-8.3; Tables8-11 & Real and decoy requests execute concurrently; only the real output is returned. This preserves the real-answer content projection, not task completion under added execution.\\
\bottomrule\end{longtable}\endgroup

